%! Constructing a Confidence Interval for a Population Proportion — Unit 6, Lesson 6.2 — Homework
\documentclass[10pt]{article}
\usepackage{apstats-article}
\usepackage{apstats-boxes}
\newcommand{\TallMath}[1]{$\displaystyle #1\rule[-1.4em]{0pt}{3.2em}$}

\begin{document}

\pageheader{Unit 6, Lesson 6.2}{Homework}
\namedateperiod

\vspace{0.15in}

\textit{For each problem, use the four-step procedure: State, Plan (verify conditions), Do (compute interval), Conclude (interpret).}

\begin{enumerate}

  \item A random sample of 400 registered voters in a city finds that 228 support a proposed
        park expansion. Construct a 95\% confidence interval for the true proportion of all
        registered voters in the city who support the expansion.

        \textbf{State:} \writelines{1}
        \textbf{Plan:}
        \begin{itemize}
          \item Random: \blank{7cm}
          \item Large Counts: \blank{7cm}
          \item 10\% rule: \blank{5cm}
        \end{itemize}
        \textbf{Do:} $\hat p = $ \blank{2cm}
        \[
          CI = \blank{1.5cm} \pm \blank{1.5cm}\sqrt{\frac{\blank{2cm}}{\blank{1.5cm}}}
             = (\blank{2.5cm},\; \blank{2.5cm})
        \]
        \textbf{Conclude:} \writelines{2}

  \item A wildlife biologist tags and releases 600 randomly selected fish from a lake.
        Later, a random sample of 80 fish from the lake is taken, and 12 are tagged.
        Construct a 99\% confidence interval for the true proportion of tagged fish in the lake.

        \textbf{State:} \writelines{1}
        \textbf{Plan:}
        \begin{itemize}
          \item Random: \blank{6cm}
          \item Large Counts: \blank{6cm}
          \item 10\% rule: \blank{5cm}
        \end{itemize}
        \textbf{Do:}
        \[
          CI = (\blank{2.5cm},\; \blank{2.5cm})
        \]
        \textbf{Conclude:} \writelines{2}

  \item A school district wants to estimate the proportion of families who would use
        an after-school tutoring program. An administrator wants a 95\% confidence interval
        with a margin of error of at most 0.05. A pilot survey suggests $\hat p \approx 0.40$.
        \begin{enumerate}[label=\alph*.]
          \item Find the minimum required sample size using the pilot estimate.
                \[n \ge \left(\frac{z^*}{m}\right)^2 \hat p(1-\hat p) = \blank{6cm} = \blank{2.5cm}\]
                Minimum $n = $ \blank{2cm}

          \item How would your answer change if you used $\hat p = 0.50$ instead? Would it be
                larger or smaller? \writelines{2}
        \end{enumerate}

\end{enumerate}

\begin{reflectionbox}
\textbf{Preview:} In Lesson 6.3 you will focus on \emph{interpreting} confidence intervals ---
what ``95\% confident'' really means for the method, and how to use an interval to evaluate a
researcher's claim about a population proportion. Think: if your CI does \emph{not} contain a
claimed value, what conclusion can you draw?
\end{reflectionbox}

\end{document}
